\documentclass[12pt]{amsart}
\usepackage{fancyhdr}
\usepackage{amsfonts}
\usepackage{amsmath,amssymb}
\usepackage[utf8]{inputenc}
\usepackage[russian]{babel}
\usepackage{amstext}
\usepackage{graphicx}
\usepackage{wrapfig}
\usepackage{soul}
\usepackage[unicode, pdftex]{hyperref}
\usepackage{ulem}
\usepackage{enumitem}
\setlist{nolistsep, itemsep=0.3cm,parsep=0pt}
\setlength{\parskip}{0.3cm}

\hypersetup
{
    colorlinks=true,
    linkcolor=blue,
    urlcolor=blue,
}

\textheight 25cm
\renewcommand{\baselinestretch}{1.22}
\voffset = -2,5cm \hoffset = -3cm \textwidth 190mm
\sloppy
\fussy

\pagestyle{fancy}
\fancyhf{} % Очищаем текущие настройки
\cfoot{\thepage} % Устанавливаем номер страницы в центре снизу
\fancyhead{}
\fancyhead[LE, LO]{\Large \it 8В} 
\fancyhead[CO,CE]{\large \it Элистинский Технический Лицей}
\fancyhead[RE,RO]{\Large \it 03.09.2024} 
\fancyfoot{}
\fancypagestyle{firststyle} 
{
\fancyhead[LO]{\large \it 8В}
\fancyhead[CO]{\Large \it Элистинский Технический Лицей}   
\fancyhead[RO]{\Large \it 03.09.2024}
\fancyfootoffset[R]{-12cm} %так можно регулировать ширину колонтитула
\fancyfoot[L]{}
\renewcommand{\footrulewidth}{0 mm} %толщина отделяющей полоски снизу
\renewcommand{\headrulewidth}{0.1 mm} %толщина отделяющей полоски сверху
\addtolength{\headheight}{5pt} % оставляем место для линейки
}

\newcommand{\problem}[2]{\noindent\llap{#1.\space}#2}

\begin{document}

\thispagestyle{firststyle}
\large

\centerline{\bf \Huge Алгебраический разнобой}

\begin{flushright}
        \scriptsize{Если я сказал не брал, значит не отдам\\ Урахаре Киске}
\end{flushright}

\problem{1}
Два различных числа $x$ и $y$ (не обязательно целых) таковы, что

$$
x^2-2000 x=y^2-2000 y .
$$


Найдите сумму чисел $x$ и $y$.

\problem{2}
Саша придумал два целых числа и перемножил их. Затем взял числа, обратные придуманным, и сложил. Полученные сумма и произведение оказались обратными числами. Найдите разность модулей Сашиных чисел.

\problem{3}
50 бизнесменов - японцы, корейцы и китайцы - сидят за круглым столом. Известно, что между каждыми двумя ближайшими японцами сидит ровно столько китайцев, сколько всего за столом корейцев. Сколько китайцев может быть за столом?


\problem{4}
Вычислите разность сумм $(1 \cdot 3+3 \cdot 5+5 \cdot 7+\ldots+99 \cdot 101)-\left(1^2+3^2+5^2+\ldots+99^2\right)$.

\problem{5}
Чему равно $\left(1+\frac{1}{2}\right)\left(1+\frac{1}{3}\right)\left(1+\frac{1}{4}\right) \ldots\left(1+\frac{1}{100}\right)$ ?


\problem{6}
Найдите значение выражения $\left(1-\frac{1}{4}\right)\left(1-\frac{1}{9}\right)\left(1-\frac{1}{16}\right) \ldots\left(1-\frac{1}{100^2}\right)$.


\problem{7}
Найдите сумму $\frac{2^2}{1 \cdot 3}+\frac{4^2}{3 \cdot 5}+\frac{6^2}{5 \cdot 7}+\ldots+\frac{100^2}{99 \cdot 101}$.

\end{document}


%\begin{wrapfigure}{r}{0.2\linewidth}
    \vspace{-7mm}
    \includegraphics[width=\linewidth]{}
\end{wrapfigure}

\begin{center}
\includegraphics[width=0.2\textwidth]{}
\end{center}2